\section{Related Work}
\label{sec:related_work}

Our system sits at the intersection of five literatures: structural
models of goal scoring, team-strength rating systems, machine
learning for match and player outcomes, fantasy-football prediction
and squad optimisation, and the economics of betting and prediction
markets. We searched arXiv, IEEE Xplore, and ResearchGate in June
2026 for work published between 2023 and 2025 on fantasy-football
prediction, World Cup outcome modelling, and market-derived features;
coverage of the most recent quarter is limited by publication lag.
This section reviews each strand in turn, summarises the practitioner
community of practice, and closes by positioning our contributions
against the closest prior art.

\subsection{Structural models of goal scoring}
\label{sec:related-structural}

The structural approach to football match prediction models goals as
Poisson arrivals at team-level rates. \citet{maher1982modelling}
introduced the independent double-Poisson model, with each side's
scoring rate derived from latent attacking and defensive strength
parameters. \citet{dixon1997modelling} refined it with a dependence
adjustment for low-scoring matches and time-decay weighting of
historical results; notably, their work was explicitly motivated by
exploiting inefficiencies in the fixed-odds betting market, which
makes it the standard football precedent for the market-versus-model
framing we adopt in Section~\ref{sec:related-markets}.
\citet{karlis2003analysis} generalised the framework to a bivariate
Poisson model with correlated goal counts and popularised the Skellam
distribution for goal differences. The independence assumption is
empirically misspecified but only mildly so:
\citet{petretta2025dependence} report a home--away goal correlation
of $-0.085$ ($p < 0.001$) across 9{,}130 matches from the five major
European leagues, and the Dixon--Coles correction redistributes
probability mass only among the four lowest scorelines, leaving
total-goals probabilities unchanged. Our structural backend
(Section~\ref{sec:methodology}) borrows the double-Poisson shape but
replaces estimated latent strengths with observable proxies (squad
price and country Elo gap), trading a small amount of fidelity for a
model that requires no training and remains interpretable end to end.

The closest structural precedent for our tournament setting is
\citet{gilch2018elo}, who combine eloratings.net Elo covariates,
Poisson regression, and Monte Carlo simulation (100{,}000 tournament
runs) to forecast the 2018 FIFA World Cup. They backtest four Poisson
variants on the 2010 and 2014 tournaments using Brier and ranked
probability scores, find a nested Poisson specification best, and
show that freezing Elo across a simulated tournament shifts stage
probabilities by up to five percentage points in favour of stronger
teams. Their architecture anticipates ours at the team level; the
per-player fantasy-points layer, which is where the difficulty of our
problem concentrates, is absent.

Textbook treatments deserve a note on lineage.
\citet{winston2022mathletics} covers least-squares team ratings,
Poisson goal models with Skellam-based outcome probabilities (citing
\citealp{karlis2003analysis}), Elo ratings, Monte Carlo tournament
simulation, and, most relevantly, formulates daily-fantasy lineup
construction as a knapsack integer program, the direct textbook
precedent for our squad optimiser. It does not, however, engage the
Maher, Dixon--Coles, or Benter lineage at all, so we cite the primary
sources directly. \citet{eager2023football} is exclusively about
American football; we cite it only for transferable methodology:
Poisson regression for count outcomes, the stability and
regression-to-the-mean behaviour of player statistics,
simulation-based projections, and open-data reproducibility practice.

\subsection{Team-strength ratings}
\label{sec:related-ratings}

Team-strength ratings supply the fixture-difficulty signal that both
our structural and learned backends consume. The Elo system
\citep{elo1978rating} originated in chess; football adaptations add
goal-margin and match-importance multipliers and retain the
calibrated scale on which a 400-point gap corresponds to ten-to-one
odds. \citet{hvattum2010using} provide the standard evaluation of
Elo-based football forecasts against bookmaker odds: Elo carries
genuine predictive information but is less sharp than the market.
The pi-rating system of \citet{constantinou2013pi} maintains separate
home and away components updated from the predicted-versus-observed
goal-difference error, diminished through
$\psi(e) = 3\log_{10}(1+e)$; it outperformed the
\citet{hvattum2010using} Elo variants and was profitable against best
bookmaker odds over five Premier League seasons, and
\citet{bunker2024soccer} confirm that pi-ratings beat Elo on
goal-difference-aware tasks. The same authors later extended this
rating-driven line to long-term team performance modelling
\citep{constantinou2017towards}. We nonetheless use Elo: it is
simpler to roll forward from open match histories, easier to falsify,
and the home--away asymmetry at the heart of pi-ratings matters less
at a World Cup played almost entirely on neutral ground. The FIFA
Men's World Ranking, the governing body's monthly official ranking,
is less responsive to recent form than Elo and serves only as a
fallback in our system (Section~\ref{sec:data}).

\subsection{Machine learning for match and player outcomes}
\label{sec:related-ml}

\citet{bunker2024soccer} demonstrate that gradient-boosted tree
models trained on engineered rating features (pi-ratings, Elo,
momentum) currently set the state of the art for club-football match
outcome prediction across multiple European leagues. For the
international tournament setting, recent work integrates team-level
and player-specific metrics with dimensionality reduction and
validates on the 2022 World Cup, finding that player attributes
combined with team composition improve outcome accuracy
\citep{worldcup2025players}; it does not address fantasy-points
prediction or the club-to-country transfer problem we tackle. A related recent line forecasts full
goal-difference distributions rather than point estimates,
% TODO cite: distributional Skellam forecasting preprint (2025)
which our player-level quantile heads capture in part. A systematic
review of machine learning in sports betting
\citep{sportsbetting2024review} identifies distribution shift between
training and live data as a recurring failure mode across the field;
the model variant we rejected for exactly this reason
(Section~\ref{sec:negative-results}) is an instance of the pattern.

At the player-regression level our problem is small-scale tabular
learning. Gradient-boosted decision trees, LightGBM
\citep{ke2017lightgbm} and XGBoost \citep{chen2016xgboost} in
particular, dominate small-to-mid tabular benchmarks, and systematic
comparisons find that they consistently beat neural networks on
tabular data at moderate scale \citep{borisov2022deeptabular,
shwartz2022tabular}. We use LightGBM for three practical reasons:
native handling of missing values in features that exist at inference
time but not in training, deterministic training under fixed seeds
(which makes feature A/B tests meaningful), and fast CPU-only
training and scoring with no model-serving runtime. We deliberately
do not use neural networks: our labelled table
(Section~\ref{sec:data}) is orders of magnitude below the data scale
at which deep models earn their keep, each row carries only five
features with no hierarchical structure for representation learning
to exploit, and a tree ensemble stored as plain text can be inspected
feature by feature.

\subsection{Fantasy-football prediction and squad optimisation}
\label{sec:related-fantasy}

Fantasy Premier League (FPL) supports a large public modelling
community, and its accumulated conventions inform our design:
position-specific models outperform pooled models, recent form is the
strongest single feature, fixture difficulty is the second, and
quantile predictions are useful for differential strategies. We adopt
all four: one model per position, form features in two of our
backends, an Elo fixture-difficulty term, and quantile heads at the
10th, 50th, and 90th percentiles.

The closest published comparable to our system is OpenFPL
\citep{groos2025openfpl}: position-specific ensembles trained on FPL
and Understat data from 2020--21 through 2023--24 and validated
prospectively on 2024--25, achieving accuracy comparable to a leading
commercial service and outperforming it on high-return players (above
two points), with models and inference code released openly. The
architectural overlap with our work is substantial: position-specific
models, multi-gameweek horizons, and open-source publication. We
differ in five respects: (i) we target an international tournament
rather than a club season; (ii) we run three explicit backends
(heuristic, structural Poisson, gradient-boosted trees) and select
the best per position on held-out validation, where their ensemble is
opaque about which member dominates where; (iii) we integrate a live
international Elo rolled from an independent match-history archive
\citep{martj42}; (iv) we couple prediction to a stage-aware
mixed-integer optimiser with transfer-hit accounting, where OpenFPL
addresses prediction accuracy only; and (v) we publish a
user-in-the-loop decision log.

On the optimisation side, \citet{mishra2025fpl} formulate starting
eleven, bench, and captain selection as mixed-integer linear programs
under FPL's budget, formation, and per-club constraints, and report
that ARIMA performs best among their point predictors. Their
optimisation formulation is close to ours; their prediction component
is simpler (a single model, no quantiles, no structural backend).
Beyond these, transformer-based sentiment analysis of player news has
been combined with statistical features to improve FPL point
prediction,
% TODO cite: sentiment-enriched FPL prediction (ResearchGate, 2025)
an idea our team-news roadmap (Section~\ref{sec:future_work}) adapts
to the international setting, and earlier baseline studies applied
linear models, decision trees, and random forests to FPL scoring.
% TODO cite: Bangaru et al. (2022), IEEE

\subsection{Prediction markets and the favourite-longshot bias}
\label{sec:related-markets}

The seminal work on combining a fundamental model with public market
signals is \citet{benter1994computer}. Benter trained a logit-based
handicapping model on horse-race data and showed that its residual,
after conditioning on the parimutuel market's implied probabilities,
still carried predictive signal: a model that beats the public must
combine its own features with the public's information rather than
replace it. Prediction markets more broadly aggregate dispersed
information into prices with well-documented forecasting accuracy
\citep{wolfers2004prediction}.

Treating market prices as direct probability estimates is
nevertheless naive. \citet{thaler1988favorite} document the
favourite-longshot bias in horse racing: casual bettors
systematically underbet favourites and overbet longshots, producing
predictable inefficiencies that sharp money arbitrages but does not
eliminate. \citet{snowberg2010explaining} test competing explanations
and attribute the bias to misperception of small probabilities rather
than risk-loving preferences. \citet{levitt2004sportsbook} shows that
sportsbooks price to exploit bettor biases rather than to balance
their books: in his sample, favourites covered the point spread in
only 49.1\% of games even though bettors overwhelmingly backed them.
\citet{croxson2014information} show that market efficiency varies
with liquidity, with in-play football prices incorporating goal news
swiftly in liquid markets while thin markets carry larger biases.
The implication for our methodology is that market prices are noisy
probability estimators whose noise mixes information aggregation
(signal) with behavioural bias (anti-signal); the appropriate
response is neither to treat them as ground truth nor to discard
them, but to combine them with a fundamental model using empirically
estimated weights, exactly as \citet{benter1994computer} did for
parimutuel racing.

We could find no published academic work that uses contemporary
prediction-market venues such as Polymarket or Kalshi as a feature
source for fantasy-football prediction, despite both venues listing
contracts directly relevant to fantasy decisions: match outcomes,
first and anytime goalscorers, and tournament top scorer. Updating
the Benter combining approach for these markets and applying it to
FIFA fantasy is, to our knowledge, new; we develop the proposal in
Section~\ref{sec:future_work}.

\subsection{Practitioner practice}
\label{sec:related-practice}

Alongside the academic literature sits a large community of practice:
editorial and analytics sites such as Fantasy Football Scout, Fantasy
Football Hub, and AllAboutFPL, tooling platforms such as Premier
Fantasy Tools and FPL Copilot, and data-science write-ups from
consultancies such as Frontier Economics.
% TODO cite: Frontier Economics FPL strategy analysis (web)
The practitioner consensus on captaincy, the highest-leverage
decision per round, is to rank by expected points, anchor on recent
form, weight fixture difficulty heavily, monitor rotation and injury
risk through pre-match press conferences, and condition differential
strategy on league position: managers who lead should match the
consensus pick, while chasing managers should take one or two
low-ownership picks, with returns diminishing beyond two. Published
booster-timing guides for the World Cup 2026 format make analogous
recommendations at the chip level.

Our system encodes much of this consensus directly (expected-points
ranking, form anchoring, an Elo fixture-difficulty term, vice-captain
as the second-highest projection), so the difference is
methodological rather than tactical: we impose held-out validation
gates before shipping a backend, compare backends systematically,
optimise under formal constraints rather than editorial tiers,
propose market integration, and keep a decision log that records
model recommendation, human override, and realised outcome. The
practitioner sites deliver excellent weekly advice for their
audience; what they do not do is measure their own cumulative
accuracy in a way that later work can build on.

\subsection{Positioning against prior art}
\label{sec:related-positioning}

Table~\ref{tab:contributions-prior-art} draws the line between prior
art and this paper's contributions.

\begin{table}[t]
\centering
\small
\begin{tabular}{@{}p{0.24\linewidth}p{0.28\linewidth}p{0.40\linewidth}@{}}
\toprule
Contribution & Closest prior art & Our delta \\
\midrule
MILP squad and captain selection &
\citet{mishra2025fpl}; knapsack formulation in
\citet{winston2022mathletics} &
Stage-aware constraints (group versus knockout), transfer-hit
accounting, three-backend prediction input \\
\addlinespace
Position-specific predictors &
\citet{groos2025openfpl}; \citet{bunker2024soccer} &
Three explicit backends with per-position selection on held-out
validation, rather than an opaque ensemble \\
\addlinespace
Club-to-country transfer &
None published &
Train on club fantasy data, infer on the World Cup; document the
distribution-shift failure modes \\
\addlinespace
Live international Elo in fantasy &
\citet{hvattum2010using} &
Elo rolled from an open match archive \citep{martj42} and
prioritised over the static FIFA ranking \\
\addlinespace
Prediction-market features &
\citet{benter1994computer} &
Polymarket and Kalshi prices proposed as a feature source for FIFA
fantasy \\
\addlinespace
Ownership-aware captaincy &
None &
Monte Carlo captain selection that models the field's likely
captain distribution \\
\addlinespace
User-in-the-loop decision log &
None &
Per-round record of model recommendation, human override, and
realised outcome \\
\bottomrule
\end{tabular}
\caption{Our contributions against the closest prior art.}
\label{tab:contributions-prior-art}
\end{table}

In short, each component of our system has a clear ancestor: the
Poisson lineage of \citet{maher1982modelling} and
\citet{dixon1997modelling}, the Elo evaluation of
\citet{hvattum2010using}, the tournament simulation of
\citet{gilch2018elo}, the fantasy prediction of
\citet{groos2025openfpl}, the fantasy optimisation of
\citet{mishra2025fpl}, and the market combining of
\citet{benter1994computer}. What has not appeared before is their
combination in a single validated pipeline for an international
tournament, run live under real transfer deadlines, with its
decisions and failures documented.
